% Source: https://artofproblemsolving.com/community/c3h3131564_2023_emcc_team_round__exeter_math_club_competition
% Stable identity: aops-topic-3131564-q05
Let the inradius be $r$. For a vertex of angle $\theta$, its distance to $I$ is $\dfrac{r}{\sin(\dfrac{\theta}{2})}$. Thus the product is $(\dfrac{2r}{\sqrt3})^2(2r)^2=\dfrac{16r^4}{3}=240$, so $r^2=3\sqrt5$. The semiperimeter is the sum of the four tangent lengths from vertices, $r(2\cot60^\circ+2\cot30^\circ)=\dfrac{8r}{\sqrt3}$. Area equals inradius times semiperimeter, giving $\boxed{8\sqrt{15}}$.
\begin{asy}
usepackage("amsmath");
size(7cm);
pair A=(-1/sqrt(3),1),B=(1/sqrt(3),1),C=(sqrt(3),-1),D=(-sqrt(3),-1);draw(A--B--C--D--cycle);draw(unitcircle,blue);dot((0,0));label("$I$",(0,0),S);label("$A$",A,NW);label("$B$",B,NE);label("$C$",C,SE);label("$D$",D,SW);
\end{asy}
